%------------------------------------------------------------------------------%
\begin{exercises}

\begin{exercise}
\begin{alphalist}
    \item Usando soma de Riemann com pontos amostrais sendo as
      extremidades direitas de cada subintervalo e $n=8$, encontre
      uma aproximação para a integral \(\int_0^4 (x^2-3x)dx\).
    \item Faça uma figura com os retângulos para ilustrar a
      aproximação da parte (a).
    \item Calculando limite da soma de Riemann e usando o
      Teorema \ref{teo:integraldefinidareescrito}
      calcule \(\int_0^4 (x^2-3x)dx\).
    \item Interprete a integral \(\int_0^4 (x^2-3x)dx\) calculada em (c)
      como uma diferença de áreas.
\end{alphalist}
\end{exercise}
%------------------------------------------------------------------------------%
\begin{exercise}
  Interprete o limite
  \(
    \lim_{n \to \infty} \sum_{i=1}^{n} \frac{\cos(x_i)}{x_i^2}
  \)
  no intervalo
  \(
    \left[\sfrac{\pi}{2}, 2\pi\right]
  \).
\begin{answer}
  \(
    \int_{\frac{\pi}{2}}^{2\pi} \frac{\cos(x)}{x^2}
  \)
\end{answer}
\end{exercise}

%------------------------------------------------------------------------------%
\begin{exercise}
  Mostre que \(\int_a^b x^2 dx = \frac{b^3-a^3}{3}\).
\end{exercise}

%---------------------------------------------------------------------------------
\begin{exercise}
  Nos exercícios a seguir, esboce o gráfico dos integrandos e calcule a
  integral interpretando em termos de áreas.
  \begin{mathlist}[2]
   \item \int_{-2}^{4} \left( \frac{x}{2}+3\right) dx
   \item \int_{-1}^{2} \vert x\vert dx
   \item \int_{-5}^{5}\sqrt{25-x^2} dx
   \item \int_0^9 \left(\frac{1}{3}x-2 \right)
  \end{mathlist}
  \begin{answer}
  \begin{mathlist}[2]
   \item 21
   \item 2,5
   \item \frac{25\pi}{2}
   \item -\frac{9}{2}
  \end{mathlist}
  \end{answer}
\end{exercise}
%----------------------------------------------------------------------------------
\begin{exercise}Sejam $f$ e $g$ funções integráveis e que
\begin{mathlist}
  \item \int_1^2 f(x)dx = -4
  \item \int_1^5 f(x)dx =  6
  \item \int_1^5 g(x)dx =  8
\end{mathlist}

Use as propriedades de integral definida para determinar
\begin{mathlist}[2]
  \item \int_2^2 f(x)dx
  \item \int_5^1 g(x)dx
  \item \int_1^2 7f(x)dx
  \item \int_2^5 f(x)dx
  \item \int_1^5 [f(x)-g(x)]dx
  \item \int_1^2 [5f(x)-2g(x)]dx
\end{mathlist}
\begin{answer}
\begin{mathlist}[2]
  \item 0
  \item -8
  \item -28
  \item 10
  \item -3
  \item -2
  \item 14
\end{mathlist}
\end{answer}
\end{exercise}

\end{exercises}
%------------------------------------------------------------------------------%
